\section{Main Body} 

The main body is the core part of the thesis, and its structure is determined by the type of work. In scientific and professional theses, the main body generally consists of the following sections: 

\begin{itemize} 
	\item current state of the problem addressed in domestic and international contexts, 
	\item objectives of the thesis, 
	\item methodology and research methods, 
	\item results, 
	\item discussion. 
\end{itemize} 

\subsection{Current State of the Problem in Domestic and International Contexts} 

In the section \textit{Current State of the Problem}, the author presents available information and knowledge related to the given topic. The sources used for this section should include current publications by both domestic and international authors. 

This part should constitute approximately 30\% of the thesis. 

The conclusion of the chapter \textit{Current State of the Problem} should provide a detailed description of a problem that has not yet been addressed within the given research area, or has been approached differently, and explain how the proposed solution could lead to improved results. 

\subsection{Objectives of the Thesis} 

The \textit{Objectives of the Thesis} section clearly, concisely, and precisely defines the subject of the research. It also includes specific sub-objectives that contribute to achieving the main objective. 

\textbf{Example of an Objective Formulation}

Based on an analysis of current methodological approaches to information system (IS) development and an analysis of approaches to the transformation of CIM into PIM within the Model Driven Architecture (MDA), the primary problem can be identified as the need to move IS development to the highest possible level of abstraction. The reason for raising the level of abstraction is the need for the most accurate possible specification of user requirements for information system design. 

Therefore, the following objective was established: 

\textbf{To develop a systematic approach to the transformation of CIM into PIM within the MDA framework, according to which business process models can be transformed into selected information system design models.} 

Achieving this objective requires dividing the solution into four partial objectives: 

\begin{enumerate} 
	\item The existence of multiple notations for the formal representation of business process models introduces the problem of selecting the most appropriate notation. Therefore, the first partial objective is: 
	
	\textbf{Selection of notations that enable the formal representation of business process models while remaining understandable to all stakeholders.} 
	
	\item The PIM level is represented by design models expressed in UML. The UML standard provides many different models for different purposes. Not all of them are suitable for use in CIM-to-PIM transformation. Therefore, the second partial objective is: 
	
	\textbf{Specification of UML models representing the PIM level.} 
	
	\item Automated or semi-automated transformations between lower MDA levels are based on the formal representation of models. Automated or semi-automated CIM-to-PIM transformation is possible only if the required formal representation of business process models is specified. Such formal representation includes both standardized graphical and semantic (textual) formalisms. Therefore, the third partial objective is: 
	
	\textbf{Determination of the required standardized formalism for automated CIM-to-PIM transformation.} 
	
	\item Identifying a suitable problem and modeling its individual components provides the basis for practical verification of the proposed transformation procedures. However, the verification will be only partial because this thesis focuses solely on the highest levels of abstraction and the transformations between them. Complete verification would require carrying out the transformations all the way to source code generation. 
	
	\textbf{The final partial objective is to develop a case study that validates the proposed approach to CIM-to-PIM transformation within the MDA framework.} 
\end{enumerate} 

\subsection{Methodology and Research Methods} 

\begin{itemize} 
	\item characteristics of the research object, 
	\item work procedures, 
	\item methods of data collection and data sources, 
	\item methods used for evaluation and interpretation of results, 
	\item statistical methods. 
\end{itemize} 

\textbf{Notes on Preparing This Chapter} 

According to Wikipedia, methodology is defined as the study of methods used in scientific work. It is a collection of procedures and methods used in a particular field. It represents a conscious and purposeful process, an organized activity, or a sequence of operations that transform initial conditions of a goal-oriented activity into its intended objective (partially or fully achieved).

Theses are part of a university's research and development activities through which scientific work is conducted. 

The development of methodology depends on: 

\begin{itemize} 
	\item the topic being addressed, 
	\item the problem-solving mechanism, 
	\item the methods used, 
	\item the verification approach. 
\end{itemize} 

A useful model may be the basic systems development life cycle consisting of the phases: analysis, design, and implementation. 

\textbf{The importance of methodology lies in its usefulness. For the researcher, it should primarily:} 

\begin{itemize} 
	\item provide a systematic and meaningful approach to solving the problem, 
	\item inspire confidence that the objectives can be achieved, 
	\item reduce the perception of scientific approaches as being overly difficult or restrictive. 
\end{itemize} 

Even if the researcher does not need or is unable to develop a formal research methodology before beginning the work, it should still be prepared retrospectively. The methodology is important for the final evaluation of results. It should be documented even if the research process was conducted in an ad hoc manner. It is necessary for assessing the validity of the solution and serves as a basis for addressing similar research problems in the future. 

According to Wikipedia, a method is defined as a recognized law transformed into a rule, a set of rules, or a system of regulatory principles. 

The presentation of methods used in a dissertation may be based on the following structure: 

\begin{itemize} 
	\item Methods of scientific cognition, including logical methods: 
	\begin{itemize} 
		\item analysis, 
		\item synthesis, 
		\item abstraction, 
		\item concretization, 
		\item induction, 
		\item deduction, 
		\item \dots 
	\end{itemize} 
		
	\item Methods of scientific investigation: 
	\begin{itemize} 
		\item empirical research (observation, measurement, experiment, \dots), 
		\item logical processing of empirical data (classification, typology, correlation, \dots), 
		\item \dots 
	\end{itemize} 
\end{itemize} 

\subsection{Results and Discussion}

The results and discussion are the most important parts of the thesis. The results (the author's own findings, opinions, or solutions to substantive problems) must be logically organized and thoroughly evaluated when presented. At the same time, all findings and observations should be discussed and compared with the results of other authors. Where appropriate, the results and discussion may be combined into a single section. Together, these sections generally constitute approximately 30--40\% of the thesis.

Depending on the nature of the work, separate chapters may be prepared for \textit{Results} and \textit{Discussion of Results}, including a detailed description of procedures followed according to the methodology. 

The \textit{Results} section should present the outcomes achieved through the conducted research.

The \textit{Discussion} section is particularly important when comparing the obtained results with those reported in other studies or alternative solutions.

\subsection{Scope of the Thesis}

The recommended length of a bachelor's thesis is 30 to 40 pages (54,000 to 72,000 characters including spaces), a master's thesis 50 to 70 pages (90,000 to 126,000 characters), a dissertation thesis 80 to 120 pages (144,000 to 216,000 characters), and a habilitation thesis up to 150 pages. Pages are counted from the introduction to the conclusion. 

According to STN ISO 690, the thesis length includes the following parts: the introduction, the main body of the text, the conclusion, the list of references, citations, and footnotes.